数值比较采用独立输入上的可行数、条件遗漏量、成功与失败耗时共同评价。主实验五种方法的可行数依次为14、14、12、17、12，不能仅以某方法较少成功样本上的目标均值判断整体优劣。消融未给出LP分量引导改善共同可行输入上目标质量的显著证据，局部搜索的目标改善则被单独记录。所有结论对应本次合成任务、实现和预算；达到时限不构成原问题无解的证明。
